\section{Open Questions}

The following five questions were left open by this replication. Each
is testable, has an explicit evidentiary basis, and comes with a
concrete next-step probe.

\begin{enumerate}
\item \textbf{Sparser / non-Hermitian matrices.} Does the hybrid HHL++
  resource reduction (fewer 2q gates, flat qubit cost) hold when the
  input matrix $A$ is sparse-but-non-Hermitian or non-trivially
  structured, rather than the small dense Hermitian matrices used in
  the paper's Table 1 and this replication?
  \emph{Basis:} the paper and this replication only test small dense
  Hermitian systems with exact-in-binary eigenvalue phases; real
  workloads (PDE discretizations, Laplacians, finance covariance) are
  sparse with irrational eigenvalues that don't align to $n$-bit
  clock precision, and non-Hermitian $A$ requires a dilation that
  adds qubits.
  \emph{Next steps:} generate a stratified suite (tridiagonal Poisson-1D
  at $n=8/16/32$; random sparse Hermitian at density $0.1$ with
  $\kappa\in\{5,10,50\}$; block-Hermitian dilations of random non-Hermitian
  systems). Threshold: hybrid claim survives iff CNOT reduction stays
  $\geq 50\%$ and fidelity stays $\geq 0.95$ across the suite.

\item \textbf{Composition with block-encoding / qubitization.} Can
  hybrid HHL++ be composed with modern block-encoding / qubitization
  primitives (e.g.\ Low--Chuang) to achieve better asymptotic scaling
  in $\kappa$ and sparsity than either technique alone?
  \emph{Basis:} semiclassical-QPE is a classical-precomputation crutch
  that flattens the QPE register; block-encoding attacks a different
  axis (LCU structure and $\kappa$ dependence). Composition behavior
  is not addressed in the paper.
  \emph{Next steps:} implement a Qiskit prototype where the hybrid
  eigenvalue-inversion block is replaced with a QSVT/qubitization
  matrix-inversion primitive; compare 2q-gate count and success
  probability vs.\ pure hybrid HHL and pure QSVT on an $8\times 8$
  sparse Hermitian instance. Metric: sub-additive resource growth
  in $(\kappa, 1/\varepsilon)$.

\item \textbf{Noise robustness on non-trapped-ion hardware.} How robust
  is the hybrid step (mid-circuit measurement + classical feed-forward
  + $\gamma$-scaling selection) to gate + measurement noise on today's
  superconducting hardware (IBM Heron / Google Willow), as opposed to
  the trapped-ion Quantinuum H-series the paper measured?
  \emph{Basis:} Quantinuum H-series has substantially better mid-circuit
  measurement fidelity than superconducting platforms
  ($\sim 99.7\%$ vs.\ $97$--$99\%$); the paper's Table~2 hardware
  fidelities (98.6\%\ at 3-bit falling to 42.6\%\ at 5-bit) are
  H-series-specific, and the hybrid classical-feedback loop can amplify
  MCM error on lower-fidelity platforms.
  \emph{Next steps:} run the paper's Sec.~4 hybrid HHL circuits on IBM
  Heron via the free Qiskit Runtime open-plan queue at 3-bit and 4-bit
  precision using dynamic circuits; report success probability and
  fidelity vs.\ noiseless statevector; compare degradation slope
  against Quantinuum Table~2.

\item \textbf{End-to-end resource accounting on a real instance.} Does
  hybrid HHL++ provide an actual end-to-end speedup on a specific
  real-world PDE or finance instance (e.g.\ 128-node Poisson-1D, 32-asset
  mean-variance portfolio) once state-preparation cost for $|b\rangle$,
  amplitude-estimation cost for extracting $x$, and repeated-sampling
  overhead from post-selection are all counted?
  \emph{Basis:} HHL's asymptotic $O(\log N)$ advantage is famously
  eroded by state-prep + readout + repetition; the paper reports
  circuit resources but not the full end-to-end budget for a
  target-accuracy solution vector.
  \emph{Next steps:} pick one concrete instance (128-node 1D Poisson,
  known analytic solution); cost out state-prep (Grover--Rudolph),
  full hybrid-HHL inversion, amplitude-estimation readout of each $x_i$
  to $\varepsilon=0.01$, and post-selection repetition factor; produce
  a break-even $(N, \varepsilon)$ plot vs.\ a classical CG solve.

\item \textbf{ML-preconditioner variant of $\gamma$-selection.} Can a
  small classical-ML preconditioner (e.g.\ a graph neural net predicting
  a good clock-register scaling $\gamma$ from $A$ directly) close the
  fidelity gap at higher clock precision (5-bit and above), where the
  paper's Table~2 shows fidelity collapsing to 42.6\%\ at 5-bit on
  real hardware?
  \emph{Basis:} the paper's $\gamma$-selection procedure (\S3) is
  hand-designed; a learned predictor could exploit matrix-structure
  patterns the hand rule misses, and the 5-bit collapse suggests
  headroom.
  \emph{Next steps:} generate $\sim$500 random small ($2\times2$--$8\times8$)
  Hermitian $A$'s with paired optimal-$\gamma$ labels from an exhaustive
  sweep; train an MLP or small GNN to predict $\gamma$ from $A$'s
  eigenvalue histogram / spectral gap; compare hybrid HHL fidelity
  using ML-$\gamma$ vs.\ paper-$\gamma$ across a held-out test set at
  $n_\mathrm{clock}\in\{3,4,5\}$. Threshold: ML wins iff mean 5-bit
  fidelity improves by $\geq 10$ percentage points.
\end{enumerate}
